\chapter{Control de Código Fuente}
\label{sec:control-codigo}

\section{Repositorio y Ramas}
\label{sec:repositorio}

El código fuente de SpeechNotes se gestiona en GitHub bajo el repositorio \url{https://github.com/gamurigm/SpeechNotes}.

\subsection{Estructura de ramas}
\label{subsec:ramas}

\begin{longtable}{lll}
\caption{Estructura de ramas del repositorio}
\label{tab:ramas}\\
\toprule
\textbf{Rama} & \textbf{Propósito} & \textbf{Protección}\\
\midrule
\endfirsthead
\toprule
\textbf{Rama} & \textbf{Propósito} & \textbf{Protección}\\
\midrule
\endhead
\bottomrule
\endfoot

\texttt{main} & Rama principal de producción & Protegida (requiere PR)\\

\texttt{planning-SQA} & Rama de trabajo SQA & Protegida (requiere PR)\\

\texttt{feat/*} & Ramas de funcionalidad & Sin protección\\

\texttt{fix/*} & Ramas de corrección & Sin protección\\

\end{longtable}

\section{Procedimiento de Cambios}
\label{sec:procedimiento-cambios}

\begin{enumerate}
    \item Crear rama \texttt{feat/xxx} o \texttt{fix/xxx} desde \texttt{main}.
    \item Implementar cambios y commits atómicos.
    \item Abrir Pull Request hacia \texttt{main}.
    \item Ejecutar pipeline CI/CD (pytest, Cypress, SonarCloud).
    \item Revisión de código por al menos 1 miembro del equipo.
    \item Merge del PR tras aprobación y verificación CI.
\end{enumerate}
